\documentclass{article}
\usepackage[UTF8]{ctex}
\usepackage{amsmath,amssymb,amsfonts}
\usepackage{enumitem}
\usepackage{geometry}
\geometry{a4paper, margin=1in}
% === 错题本排版优化（不参与题目标识） ===
\setlength{\parindent}{0pt}
\setlength{\parskip}{0.35em plus 0.10em minus 0.08em}
\setlength{\emergencystretch}{3em}
\setlist{itemsep=0.25em, topsep=0.40em, parsep=0pt}
\allowdisplaybreaks[2]
% === 错题本排版优化结束 ===
\begin{document}
\section*{西北大学2024年数学分析真题}
\section*{176．西北大学2024年数学分析试卷}
\begin{itemize}
  \item[一．] （20 分）求下列极限。
  \item[（1）] \(\lim _{x \rightarrow 0}\left(\frac{2+e^{\frac{1}{x}}}{1+e^{\frac{4}{x}}}+\frac{\sin x}{|x|}\right)\) ．
  \item[（2）] \(\lim _{x \rightarrow 0} \frac{1}{x^{3}}\left[\left(\frac{2023+\cos x}{2024}\right)^{x}-1\right]\) ．
  \item[二．] （20分）设 \(x_{1}>0\) ，作迭代序列
\[
x_{n+1}=\ln \left(1+x_{n}\right), n=1,2, \cdots .
\]
  \item[（1）] 证明： \(\lim _{n \rightarrow \infty} x_{n}\) 存在，并求其值．
  \item[（2）] 求 \(\lim _{n \rightarrow \infty} \frac{x_{n} x_{n+1}}{x_{n}-x_{n+1}}\) ．
\end{itemize}

\subsection{题目1}
三．（15 分）设 \(f(x)\) 在 \([-1,1]\) 上有连续的二阶导数，且 \(f(0)=0\) ．

\begin{itemize}
  \item[（1）] 写出 \(f(x)\) 的带拉格朗日型余项的一阶麦克劳林公式。
  \item[（2）] 证明：至少存在一点 \(\eta \in[-1,1]\) ，使得 \(f^{\prime \prime}(\eta)=3 \int_{-1}^{1} f(x) \mathrm{d} x\) ．
\end{itemize}

\subsection{题目2}
四．（15 分）已知 \(z=f(x, y)\) 的全微分为 \(\mathrm{d} z=2 x \mathrm{~d} x-2 y \mathrm{~d} y\) ；且 \(f(1,1)=2\) ，求 \(f(x, y)\) 在有界闭区域 \(D: x^{2}+\frac{y^{2}}{4} \leq 1\) 上的最大值与最小值．

\subsection{题目3}
五．（15 分）设 \(\Omega\) 为半球 \(\left\{(x, y, z) \mid x^{2}+y^{2}+z^{2} \leq 1, z \geq 0\right\}\) ，计算

\[
I=\iiint_{\Omega}(\sqrt{2} x+z)^{2} \mathrm{~d} x \mathrm{~d} y \mathrm{~d} z .
\]

\subsection{题目4}
六．（15 分）设 \(P(x, y)=\frac{x\left(x^{2}+y^{2}\right)^{k}}{y}, Q(x, y)=-\frac{x^{2}\left(x^{2}+y^{2}\right)^{k}}{y^{2}}\) ．

\begin{itemize}
  \item[（1）] 若积分 \(\int_{L} P \mathrm{~d} x+Q \mathrm{~d} y\) 在 \(D=\{(x, y) \mid y>0\}\) 上与路径无关，求 \(k\) 的值。
  \item[（2）] 对第（1）的 \(k\) ，在 \(D\) 内求函数 \(u(x, y)\) ，使得 \(\mathrm{d} u=P \mathrm{~d} x+Q \mathrm{~d} y\) ，并计算 \(I=\int_{(1,2)}^{(3,4)} P \mathrm{~d} x+Q \mathrm{~d} y\) ．
\end{itemize}

\subsection{题目5}
七．（15 分）计算第二类曲面积分

\[
\iint_{\sum} \mathrm{d} y \mathrm{~d} z+z \mathrm{~d} z \mathrm{~d} x+\frac{e^{x}}{\sqrt{x^{2}+y^{2}}} \mathrm{~d} x \mathrm{~d} y .
\]

其中 \(\sum\) 是曲面 \(z=\sqrt{x^{2}+y^{2}}\) 介于两平面 \(z=1\) 和 \(z=2\) 之间的部分，取下侧．

\subsection{题目6}
八．（15 分）求函数项级数 \(\sum_{n=1}^{\infty}\left(x+\frac{1}{n}\right)^{n}\) 的收敛域，并讨论和函数的连续性．

\subsection{题目7}
九．（20 分）设 \(\varphi(x)=\int_{0}^{+\infty} \frac{t}{2+t^{x}} \mathrm{~d} t(x \geq 3)\) ．令

\[
f_{n}(x)=\sum_{k=0}^{n-1} \frac{1}{n} \varphi\left(x+\frac{k}{n}\right), n=1,2, \cdots .
\]

证明：函数列 \(\left\{f_{n}(x)\right\}\) 在 \([3, A]\) 上一致收敛，其中 \(A>3\) ．
\end{document}
